% conclusion.tex — Conclusion
\section{Conclusion}
\label{sec:conclusion}

We presented Voice \& Vision Assistant (VVA), an open-source, multimodal real-time accessibility system that fuses twelve perception and interaction capabilities---object detection, monocular depth estimation, edge-aware segmentation, OCR, braille reading, QR/AR scanning, face recognition, sound localization, action recognition, visual question answering, retrieval-augmented memory, and natural speech---into a single spoken assistant for blind and visually impaired users.

Three architectural decisions distinguish VVA from prior work. First, the layered architecture with import-linter enforcement at CI time guarantees that no perception module accidentally depends on infrastructure or application code, making each component independently testable (429+ tests and counting). Second, the parallel perception pipeline with per-stage timeouts and multi-tier fallbacks keeps the system responsive within a 500\,ms end-to-end budget on consumer GPU hardware, even when individual components degrade---a practical necessity that many academic systems overlook. Third, the privacy-first design (local FAISS indexing, opt-in consent, GDPR erasure/portability endpoints) makes VVA deployable in healthcare and education settings where cloud-resident personal data is unacceptable.

We demonstrated that the perception stack---YOLOv8n detection, MiDaS~v2.1 depth, edge-aware segmentation, and spatial fusion---processes frames at a median latency of 62\,ms on RTX~4060, well within the 300\,ms perception budget. The priority-weighted risk scoring formula (Eq.~\ref{eq:risk}) produces actionable micro-navigation cues that convey the most important hazard in under ten words.

Significant work remains. The COCO class set is too narrow for assistive use; metric depth is still approximated; and formal user studies are essential before any real-world deployment claim can be substantiated. We release the full codebase, model checkpoints, pipeline configurations, and test suite as open source to enable community replication, extension, and rigorous evaluation.

Ultimately, the goal of VVA is not to replace the white cane or the guide dog, but to add an intelligent layer of environmental awareness that speaks when something matters---and stays silent when the path is clear.
